\section{引言与问题界定}
\label{sec:intro}

在电力系统运行与电力现货市场中，风电功率预测的实际运行价值根本上取决于日前发电计划申报与实时平衡调度中的非对称决策损失 \cite{holttinen2021,weron2014}。在典型的两部制平衡结算规则下，预测不足导致的电网短缺惩罚与多发导致的弃电回收价格存在显著差异。在这种非对称决策工况下，均方根误差 (RMSE) 或平均绝对误差 (MAE) 等传统对称回归指标无法代表下行经济风险与电网备用安全敞口 \cite{lipton2019}。为提升模型泛化性能并防止产生违背机组运行特性的反常预测，大量研究倡导将气动极限、经验功率曲线或运动学方程直接融入深度神经网络结构与损失函数中 \cite{raissi2019physics,karniadakis2021physics,sheng2025annreview,pct2026}。这一研究范式普遍建立在一个核心前提之上：所报告的预测误差下降，能够直接归因于物理约束中编码的有效领域知识。

然而，现行的文献报告惯例、开源代码赤字与消融实验设计掩盖了这一归因基础。首先，现有文献几乎普遍报告积极收益，极少探讨物理先验的失效工况或负面权衡。其次，公开可复现的源代码极其匮乏，导致算法的底层计算图执行往往缺乏独立检验 \cite{kapoor2023,pineau2021}。再次，现有消融实验普遍通过直接剔除多层神经网络子模块（如特征编码通道或注意力模块）来评估物理约束，而非在保持网络容量不变的前提下隔离物理惩罚权重 $\lambda$。这种粗粒度消融不可避免地将网络参数容量与物理归纳偏置混为一谈。因此，现有文献留下了一个核心追问：物理约束带来的收益，究竟来自有效的先验信息，还是其他因素？这些收益在什么条件下会消失或逆转？

为了从经验上检验物理约束收益的归因依据，我们在三座英国商业风电场（Kelmarsh、Hill of Towie 与 Penmanshiel）开展了受控实证审计。在保持预测网络骨干 (DLinear) 与验证集后验分位数偏移校准完全一致的前提下，我们引入形式与权重匹配的伪打乱先验作为阴性对照，系统扰动先验参数并改变模型训练窗口。受控实验表明，物理先验的效益高度依赖于具体数据条件与约束设置：
\begin{enumerate}[leftmargin=*]
\item \textbf{先验特异性与误设脆弱性}：在 Penmanshiel 样本受限工况下（$N_{\text{train}}=1{,}583$），施加强功率曲线先验使归一化非对称决策损失降低 1.22\%。然而，在保持损失形式、模型架构与正则化权重完全相同的条件下，打乱风速与功率的物理对应关系会导致决策损失激增 68.1\%。这一鲜明对比证实：虽然有效领域知识在数据匮乏时能提供可测量的归纳保护，但强正则化下模型对先验误设表现出极高的脆弱性。
\item \textbf{信息预算与样本机制}：在不同训练样本预算的评估中，固定历史经验先验在较短训练窗口下（如 $N=500$，约对应 21 天近期样本）展现出边际优势。然而，由于该经验先验是在全年完整历史数据（7 至 10 个月）上拟合所得，它实际为模型注入了显著的额外历史信息预算。随着模型自身训练样本扩大至 $N=5{,}000$，先验优势逐渐消退并发生轻微反转，刚性的历史曲线反而限制了模型适应局部动态变化的能力。
\item \textbf{实现偏离与市场指标脱钩}：在软件执行层面，对开源计算图的核查发现物理约束 Transformer (PCT) 的 Jensen--Shannon (JS) 散度损失函数存在反向参数传递，引发了非对称的反向梯度缩放。在下游市场评估中，对结算机制的数理分析表明，孤立的未对冲平衡现金流与预测能力存在结构性脱钩：退化零申报基准（$q \equiv 0$）无需任何预测能力即可机械获取巨额名义现金流。
\end{enumerate}

\subsection{主要研究贡献}
本文从文献报告规范、底层软件执行与不平衡结算及决策风险诊断三个相互支撑的维度，对物理约束风电预测展开系统实证审计：
\begin{itemize}[leftmargin=*]
\item \textbf{文献报告规范与代码可用性范围界定审计}：在 PRISMA-ScR 指导下，我们对 46 项当代代表性研究（2023--2026）展开全生命周期范围界定审计，量化了失效边界报告的缺失（0/46，0.0\%）、开源代码的严重赤字（44/46，95.7\%）以及模块级结构消融的普遍混淆（23/46，50.0\%）。
\item \textbf{骨干对齐的前瞻性受控消融基准}：通过在 10 卡专用物理 GPU 集群上跨三座风电场（各站 3--4 个独立随机种子）将真实先验与尺度对齐的伪打乱阴性对照进行成对消融，利用形式与权重对齐的阴性对照区分领域特定先验效应与非特异性数值平滑，刻画了物理约束收益在样本规模、先验扰动与气象代理时效错配下的实证边界。
\item \textbf{底层计算图核查与结算脱钩诊断}：我们揭示了开源物理约束软件中存在的实现偏离，并通过数理与实证分析阐明了为何孤立的单边平衡现金流不能作为预测质量的评价指标，确立了以归一化非对称调度损失为核心的决策理论审计框架。
\end{itemize}

\paragraph{研究范畴与方法边界}
本文确立四项明确的实证审计范畴：首先，文献范围界定审计刻画纳入的 46 项当代研究（38 篇期刊论文与 8 篇预印本，2023--2026）的报告与消融规范，反映该范围界定样本内的报道现状而非断言全领域通用属性；其次，软件审计基于开源 PCT 代码库（commit \texttt{65941ec}）分析其静态计算图实现保真度；第三，前瞻性受控消融实验在三座商业风电场上以对齐的 DLinear 为基准刻画边界行为，并辅以 PatchTST 作为结构敏感性探针；第四，下游市场评估基于英国平衡机制两部制非对称短缺清算规则开展不平衡结算与决策风险诊断，揭示评价指标的病态反常而非承诺全组合商业清算。上述四个维度共同构成扎实的实证证据链。
